\documentclass[11pt]{article}

\usepackage[a4paper,margin=27mm]{geometry}
\usepackage{amsmath,amssymb,amsthm,mathtools,bm}
\usepackage{booktabs}
\usepackage{microtype}
\usepackage[hidelinks]{hyperref}

\newtheorem{theorem}{Theorem}
\newtheorem{lemma}[theorem]{Lemma}
\newtheorem{proposition}[theorem]{Proposition}
\newtheorem{remark}[theorem]{Remark}
\newcommand{\C}{\mathbb C}
\newcommand{\Q}{\mathbb Q}
\newcommand{\ket}[1]{\lvert #1\rangle}
\newcommand{\bra}[1]{\langle #1\rvert}
\newcommand{\opnorm}[1]{\left\lVert #1\right\rVert}
\newcommand{\Tq}{T_q}

\title{Exact all-parameter spectral-gap certificates for\\
deformed-EPR Hamiltonians on small subdivided stars}
\author{Xiaojian Zeng\\
\small PBC School of Finance, Tsinghua University, Beijing 100190, China\\
\small \texttt{xiaojianzeng@foxmail.com}}
\date{}

\begin{document}
\maketitle

\begin{abstract}
We study the coefficient-one deformed-EPR Hamiltonian
\(H_s=-\sum_{e}(\ket{00}+s\ket{11})(\bra{00}+s\bra{11})_e\)
on finite trees, where \(0<s<1\) and the rank-one edge term is not
normalized.  A recent conjecture gives the lower bound
\(\Delta(H_s)\geq 1-s^2\), with equality on stars.  We give an exact
finite-family result for the subdivided star \(T_q\), obtained from a star
by subdividing one edge: for every \(q=1,\ldots,6\) and every real
\(0<s<1\),
\[
  \Delta(T_q,s)>1-s^2.
\]
The case \(q=1\), the four-vertex path, is proved by an explicit quartic
root-ordering argument.  For \(q=2,\ldots,6\), permutation symmetry and a
conserved staggered occupation reduce the full Hamiltonian to real
symmetric blocks of dimension at most eight.  Exact characteristic
polynomials and crossing resultants over \(\Q[t,x]\), with \(t=s^2\),
reduce a possible failure to 171 algebraic parameter neighborhoods.
Rational root isolation and Weyl-continuity enclosures prove a positive
margin on every neighborhood, while an exact reference sign fixes each
remaining component.  The result excludes these six non-star trees as
additional equality cases.  It does not prove the conjecture for arbitrary
trees or for \(q\geq7\).
\end{abstract}

\section{Introduction}

Wong, Bravyi, Gosset and Liu introduced a family of two-local Hamiltonians
whose edge interactions favor a deformed EPR vector and connected its
spectral properties to Lee--Yang tensors~\cite{wong2026}.  In the
coefficient-one convention used here, the Hamiltonian on a finite graph
\(G=(V,E)\) is
\begin{equation}
 H_s(G)=-\sum_{\{i,j\}\in E}
 \bigl(\ket{00}+s\ket{11}\bigr)
 \bigl(\bra{00}+s\bra{11}\bigr)_{ij},
 \qquad 0\leq s<1.                                      \label{eq:H}
\end{equation}
This model lies in the Suzuki--Fisher class~\cite{suzuki1971}.  The source
proved uniqueness and even parity of the ground state and found the exact
gap \(1-s^2\) on every star with at least three vertices.  It then
conjectured the lower bound
\begin{equation}
       \Delta(G,s)\geq 1-s^2                            \label{eq:conj}
\end{equation}
for every finite connected graph.  Its evidence included exact
diagonalization at 100 uniformly spaced parameter values for all trees with
at most 18 vertices and all connected graphs with at most 10 vertices, as
well as a mixed exact-diagonalization and matrix-product-state study of
paths~\cite{wong2026}.  Those finite grids are strong evidence, but do not
order the spectrum over an uncountable parameter interval.

Two subsequent Lee--Yang gap results give broad but differently normalized
bounds.  Bravyi, Gosset, Liu and Wong proved a gap of at least twice the
minimum longitudinal field for the Suzuki--Fisher family~\cite{bravyi2026}.
Using the Pauli expansion of~\eqref{eq:H}, this gives
\(\Delta\geq \tfrac12(1-s^2)\delta(G)\), where \(\delta(G)\) is the
minimum degree.  It therefore supplies only half of~\eqref{eq:conj} on a
tree.  Rayudu and Takahashi proved another uniform-field Lee--Yang gap
bound~\cite{rayudu2026}.  Neither result provides the sharp all-parameter
ordering established below for the specified finite family.

Our graph \(T_q\) has ordinary leaves \(L_1,\ldots,L_q\), a centre \(C\),
and two further vertices \(A,B\), with edges
\begin{equation}
  CL_i\quad(1\leq i\leq q),\qquad CA,\qquad AB.          \label{eq:Tq}
\end{equation}
Thus \(T_q\) has \(q+3\) vertices and \(q+2\) edges; it is obtained by
subdividing one edge of the star \(K_{1,q+1}\).  Every \(T_q\) is a
non-star, including \(T_1=P_4\).  The main result is deliberately finite.

\begin{theorem}[Finite subdivided-star result]             \label{thm:main}
For each integer \(q\in\{1,2,3,4,5,6\}\) and every real
\(0<s<1\), the Hamiltonian~\eqref{eq:H} on \(T_q\) satisfies
\begin{equation}
             \Delta(T_q,s)>1-s^2.                        \label{eq:main}
\end{equation}
\end{theorem}

The restriction in \(q\) is part of the theorem, not a finite-size proxy
for an asymptotic claim.  The proof is analytic for \(q=1\) and exact
computer-assisted for \(2\leq q\leq6\).  It controls both the lowest odd
state and the second even state.  This is essential: controlling only the
odd excitation would assume a separate conjecture made in
Ref.~\cite{wong2026}.

\section{Hamiltonian, normalization and gap channels}

Put \(t=s^2\) and work with the positive semidefinite operator
\begin{equation}
 K_s=-H_s=\sum_{e\in E}\ket{\phi_s}\!\bra{\phi_s}_e,
 \qquad \ket{\phi_s}=\ket{00}+s\ket{11}.              \label{eq:K}
\end{equation}
In the ordered basis \((\ket{00},\ket{01},\ket{10},\ket{11})\), one edge
term is
\begin{equation}
 \ket{\phi_s}\!\bra{\phi_s}=
 \begin{pmatrix}
 1&0&0&s\\0&0&0&0\\0&0&0&0\\s&0&0&s^2
 \end{pmatrix}.                                       \label{eq:local}
\end{equation}
It has rank one and sole nonzero eigenvalue \(1+s^2\).  In particular,
it is not a projector in the convention of~\eqref{eq:H}.  Energies and
gaps below are dimensionless, with coefficient one on every edge.

The parity \(P=Z^{\otimes |V|}\) commutes with \(K_s\).  Let \(\mu_e\)
and \(\mu_o\) be the largest eigenvalues in the even and odd sectors,
and let \(\nu_e\) be the second-largest even eigenvalue.  The ground-state
uniqueness result of Ref.~\cite{wong2026} implies that the top eigenvalue
of \(K_s\) is \(\mu_e\).  Reversing the order of the spectrum when
passing from \(K_s\) to \(H_s=-K_s\) gives
\begin{equation}
 \Delta(T_q,s)=\mu_e-\max\{\mu_o,\nu_e\}.              \label{eq:gap}
\end{equation}
It is therefore enough, and necessary for our method, to prove positivity
of both margins
\begin{equation}
 M_o(t)=\mu_e-\mu_o-(1-t),\qquad
 M_e(t)=\mu_e-\nu_e-(1-t).                            \label{eq:margins}
\end{equation}

The endpoint behavior illustrates why the strict statement uses the open
interval.  At \(s=0\), \(K_0\) is diagonal and counts zero--zero edges;
every finite tree has \(\Delta=1\).  At \(s=1\) the target bound vanishes.
Neither endpoint is included in Theorem~\ref{thm:main}.

\section{Spin and staggered-charge reduction}              \label{sec:blocks}

Permutations of the \(q\) ordinary leaves commute with \(K_s\).  Their
Hilbert space decomposes into spin
\begin{equation}
 j=\frac q2-r,\qquad 0\leq r\leq\left\lfloor\frac q2\right\rfloor,
 \qquad
 d_{q,r}=\binom qr-\binom q{r-1},                       \label{eq:spin}
\end{equation}
where \(d_{q,r}\) is the multiplicity and \(\binom q{-1}=0\).  In one
copy of a spin sector, use the orthonormal basis
\(\ket{c,a,b;k,r}\), where \(c,a,b\in\{0,1\}\) are the bits at
\(C,A,B\), and \(r\leq k\leq q-r\) is the physical Hamming weight of
the ordinary leaves.  The diagonal contribution from the \(q\) ordinary
edges is
\begin{equation}
 q-k\quad(c=0),\qquad tk\quad(c=1),                    \label{eq:diag}
\end{equation}
and their nonzero off-diagonal matrix element is
\begin{equation}
 \ket{0,a,b;k,r}\longleftrightarrow\ket{1,a,b;k+1,r},
 \qquad s\sqrt{(q-r-k)(k-r+1)}.                        \label{eq:offdiag}
\end{equation}
The edges \(CA\) and \(AB\) add the same two-state action as
Eq.~\eqref{eq:local} on the corresponding core bits.

There is a further reduction that uses bipartiteness.  Assign positive
sign to \(L_i,A\), negative sign to \(C,B\), and define
\begin{equation}
                 Q=k+a-c-b.                            \label{eq:charge}
\end{equation}
Every off-diagonal transition flips \(00\leftrightarrow11\) across the
bipartition, so \(Q\) is conserved.  Moreover,
\((k+a+c+b)-Q=2(c+b)\), hence parity equals \(Q\bmod2\).  For fixed
\((r,Q)\), Eq.~\eqref{eq:charge} determines
\(k=Q-a+c+b\) from three core bits.  Each real symmetric block
\(B_{q,r,Q}(t)\) therefore has dimension at most eight, independently of
\(q\).

The complete parity characteristic polynomial factors as
\begin{equation}
 \chi_p(x,t)=
 \prod_{r=0}^{\lfloor q/2\rfloor}
 \prod_{Q\equiv p\ ({\rm mod}\ 2)}
 \det\!\left[xI-B_{q,r,Q}(t)\right]^{d_{q,r}}.          \label{eq:factor}
\end{equation}
The dimension identity
\begin{equation}
 \sum_r d_{q,r}\sum_Q\dim B_{q,r,Q}=2^{q+3}            \label{eq:dimension}
\end{equation}
is checked symbolically for every certified \(q\), so no permutation or
parity channel is omitted.  Since each parity matrix is entrywise
nonnegative, a top eigenvector can be chosen nonnegative and averaged over
leaf permutations without vanishing.  Thus the top eigenvalue in either
parity lies in \(r=0\).  At \(s=0\), the unique top state is the all-zero
state in \((r,Q)=(0,0)\).  Ground-state simplicity for \(0\leq s<1\)
prevents this root from changing invariant sectors.  Consequently the
physical ground root is always the largest root of the \((0,0)\) block.
For \(q\geq2\), after an identically zero factor is removed, its
characteristic polynomial has degree six.

\section{Analytic proof for \texorpdfstring{\(T_1=P_4\)}{T1 = P4}}
\label{sec:p4}

For the four-vertex path, exact construction of the parity blocks gives
\begin{equation}
 \chi_e(x)=x^3(x-1-t)Q_t(x),                            \label{eq:chie}
\end{equation}
where
\begin{align}
Q_t(x)={}&x^4-5(1+t)x^3+(7+18t+7t^2)x^2 \nonumber\\
 &-3(1+t)(1+4t+t^2)x+2t(1+4t+t^2),                    \label{eq:quartic}
\end{align}
and
\begin{equation}
 \chi_o(x)=x^2(x-1-t)^2
 \left[x^2-2(1+t)x+2t\right]^2.                       \label{eq:chio}
\end{equation}
The largest odd eigenvalue is
\begin{equation}
 b(t)=1+t+r,\qquad r=\sqrt{1+t^2}.                   \label{eq:b}
\end{equation}
All roots of \(Q_t\) are real because it is a factor of the
characteristic polynomial of a real symmetric invariant block.  At
\(t=0\), \(Q_0=x(x-1)^2(x-3)\), so exactly one root of \(Q_0\) lies
above \(b(0)=2\).  Reduction modulo \(r^2=1+t^2\) yields the exact identity
\begin{equation}
 Q_t(b(t))=-(1-t)^2(1+t)b(t)<0\qquad(0<t<1).           \label{eq:qb}
\end{equation}
No root can cross the continuous threshold \(b(t)\) without making the
left side zero.  Hence exactly one even root remains above \(b(t)\), and
the first excitation is odd throughout the interval.

At the target-shifted odd root, the second exact identity is
\begin{equation}
 Q_t\bigl(b(t)+1-t\bigr)=Q_t(2+r)
 =-8t(1-t)\bigl[(1-t)r+t^2-t+1\bigr]<0.               \label{eq:qshift}
\end{equation}
Since \(Q_t\) is monic, its largest root is strictly larger than
\(b(t)+1-t\).  That root is \(\mu_e\), so
\begin{equation}
       \Delta(P_4,s)=\mu_e-b(t)>1-t=1-s^2.            \label{eq:p4done}
\end{equation}
This proves Theorem~\ref{thm:main} for \(q=1\).

\section{Exact crossing certificates for
\texorpdfstring{\(2\leq q\leq6\)}{2 <= q <= 6}}
                                                                  \label{sec:cert}

This section describes the finite, independently checkable certificate
used for the remaining cases.  Let \(P_q(x,t)\) be the nonzero degree-six
factor containing the ground root, and let \(F_{q,r,Q}(x,t)\) run over
all nonzero charge--spin factors in~\eqref{eq:factor}.  The ground factor
itself is included when comparing its lower roots.  All these polynomials
are computed over \(\Q[t,x]\).

If either margin in~\eqref{eq:margins} vanished at an interior parameter,
then for some physical competitor \(\lambda\),
\begin{equation}
 P_q(\mu_e,t)=0,
 \qquad
 F_{q,r,Q}\bigl(\mu_e-(1-t),t\bigr)=0.                \label{eq:crossing}
\end{equation}
Thus \(t\) must be a zero of the exact resultant
\begin{equation}
 R_{q,r,Q}(t)=\operatorname{Res}_x
 \left(P_q(x,t),F_{q,r,Q}(x-(1-t),t)\right).           \label{eq:res}
\end{equation}
For every listed block, this resultant is nonzero.  Powers of \(t\) and
\(t-1\) are removed because the theorem concerns the open interval, and
the remaining univariate polynomial is replaced by its square-free part.
Exact real-root isolation then returns every root in \((0,1)\).  Across
\(q=2,\ldots,6\), retaining duplicates between different resultants,
there are 171 candidate intervals.

At the rational midpoint \(t_0\) of each isolating interval \([a,b]\),
all real roots of all physical block polynomials are isolated over
\(\Q\).  Taking a rational lower endpoint for the ground root and rational
upper endpoints for every competing root gives an exact lower bound
\(m_0\) on
\begin{equation}
 g_q(t_0)=\mu_e(t_0)-\max\{\mu_o(t_0),\nu_e(t_0)\}-(1-t_0).
                                                                    \label{eq:g}
\end{equation}
The pointwise bound must also control every parameter inside the isolating
interval.  Differentiating one edge term in the
\((\ket{00},\ket{11})\) subspace gives
\begin{equation}
 \opnorm{\frac{d}{ds}\ket{\phi_s}\!\bra{\phi_s}}
 =s+\sqrt{1+s^2}<\frac52.                              \label{eq:deriv}
\end{equation}
With \(q+2\) edges, Weyl's inequality therefore makes every ordered
eigenvalue \(L_q\)-Lipschitz in \(s\), where
\begin{equation}
                    L_q=\frac52(q+2).                 \label{eq:L}
\end{equation}
The maximum of the competitors has the same Lipschitz constant.  If
\(t_0=(a+b)/2\), then
\begin{equation}
 |\sqrt t-\sqrt{t_0}|
 =\frac{|t-t_0|}{\sqrt t+\sqrt{t_0}}
 \leq \frac{b-a}{4a}.
\end{equation}
Combining the variation of two eigenvalues with that of \(1-t\) gives the
rational enclosure
\begin{equation}
 |g_q(t)-g_q(t_0)|
 \leq \frac{b-a}{2}\left(\frac{L_q}{a}+1\right).       \label{eq:weylbound}
\end{equation}
For every one of the 171 candidate intervals, the exact lower bound
\(m_0\) is strictly larger than the right side of
Eq.~\eqref{eq:weylbound}.  Table~\ref{tab:cert} summarizes the inventory.
The displayed decimals are readability aids; every comparison is stored
and checked as a rational inequality.

\begin{table}[ht]
\centering
\caption{Exact certificate summary.  ``Blocks'' counts nonzero
charge--spin factors.  Candidate intervals retain duplicates across
resultants.  The reference margin is at \(t=1/4\).}
\label{tab:cert}
\begin{tabular}{@{}rrrrrr@{}}
\toprule
\(q\)&\(n\)&blocks&intervals&reference margin&minimum slack\\
\midrule
2&5& 8&18&0.1072347684&0.0196245177\\
3&6&10&21&0.0668105392&0.0079834888\\
4&7&16&33&0.0486643413&0.0041072985\\
5&8&19&41&0.0383207710&0.0024037224\\
6&9&26&58&0.0316288710&0.0010475231\\
\bottomrule
\end{tabular}
\end{table}

Exact root isolation also gives \(g_q(1/4)>0\) for every \(q\) in the
table.  If \(g_q\) vanished anywhere in \((0,1)\),
Eq.~\eqref{eq:res} would place that point in a certified neighborhood on
which \(g_q>0\), a contradiction.  Continuity and the positive reference
sign therefore prove \(g_q(t)>0\) throughout \((0,1)\).  Identically zero
eigenvalue factors cannot be competitors: the all-zero trial state has
Rayleigh quotient \(q+2\), so \(\mu_e-0>1-t\).  This completes the proof
of Theorem~\ref{thm:main}.

\subsection{Independent representation check}

The exact generator uses fixed-size collective-weight matrices and checks
Eq.~\eqref{eq:dimension}.  A separate checker constructs both parity
matrices directly in the full computational basis from the local action
\(00\leftrightarrow11\).  It verifies every saved rational neighborhood
inequality and compares the full-basis spectrum at the exact reference
point and the least-slack midpoint for each \(q\).  These floating-point
diagonalizations corroborate the representation but are not used to prove
the all-parameter result.  The proof rests on the exact polynomials,
resultants, rational isolators and enclosure~\eqref{eq:weylbound}.  The
symbolic calculations use SymPy~\cite{meurer2017}; the software environment
and commands are specified in the accompanying reproducibility notes.

\section{Physical interpretation and limitations}

The star value \(1-s^2\) is the known tight benchmark for the conjectured
graph-independent lower bound~\cite{wong2026}.  Theorem~\ref{thm:main}
shows that subdividing one star edge does not create another equality case
for the first six members of this family: on four through nine qubits the
finite-volume gap is pointwise strictly larger for every \(0<s<1\).  In an
adiabatic interpolation that stays inside this finite model, these graphs
therefore have a larger instantaneous gap than the star benchmark at the
same \(s\).

Several qualifications are essential.  First, the strict excess is not
uniform on the closed parameter interval; it can vanish as an endpoint is
approached.  Second, the cutoff \(q\leq6\) is arbitrary and the certificate
does not establish the whole subdivided-star family.  Third, no claim is
made for an arbitrary tree, a graph with cycles, a weighted interaction, a
normalized projector, a phase-deformed edge term, or the thermodynamic
limit.  Finally, a finite-tree result alone does not supply a new algorithm
for arbitrary bipartite Quantum Max-Cut.  Broader algorithmic consequences
now follow from the general result of Ref.~\cite{bravyi2026}, by a route
different from this finite-family certificate.

The block reduction suggests a sharper next step.  Nonsymmetric leaf-spin
sectors reduce to smaller symmetric problems with a positive but nonscalar
field at \(C\).  A branch-selective Schur-complement or interlacing argument
could compare the \((r,Q)=(0,0)\) Perron root directly with the leading
\((0,1)\), \((0,2)\), and recursive \(r>0\) blocks.  Such an argument would
avoid the irrelevant crossings that cause the resultant inventory to grow,
and could either prove the result for every \(q\) or isolate the first exact
failure.  That extension remains open.

\section{Conclusion}

We converted sampled evidence for six small non-star trees into an exact
continuous-parameter statement.  The combination of a fixed-size
spin--charge decomposition, exact crossing resultants, rational root
isolation and Weyl enclosures proves
\(\Delta(T_q,s)>1-s^2\) for \(q=1,\ldots,6\) and all real \(0<s<1\).
The method controls all excitation sectors and distinguishes its exact
certificate from a finite grid scan.  The universal tree bound and its
complete equality classification remain unresolved.

\section*{Data and code availability}

The Python generator, exact JSON certificate, independent full-basis
checker, environment specification and regeneration instructions are
included as supplementary files with this draft.  The certificate uses
open rational formats and does not require proprietary data.  A permanent
public repository and identifier remain to be selected by the author before
submission.

\section*{Generative-AI disclosure}

OpenAI Codex using the \texttt{gpt-5.6-sol} model assisted with problem
development, literature-search support, symbolic derivations, code
generation and execution, same-machine checking, and manuscript drafting.
This local draft and its calculations require critical review and final
approval by the named author before submission.  The AI system is not an
author.

\section*{Declaration of interests}

The declaration of interests is pending confirmation by the named author.

\bibliographystyle{unsrt}
\bibliography{references}

\end{document}
